\documentclass{article}
\usepackage[a4paper, total={6.5in, 10in}]{geometry}
\setlength{\parindent}{0pt}
\usepackage{tcolorbox}
\tcbset{colback=white!94!black, colframe=black!60!white, coltitle = white, boxrule=0.8pt, arc=2mm}
\newtcolorbox{boxs}[2][]{title= #2,#1}
\usepackage{datetime2}
\usepackage[british]{babel}
\usepackage{amsfonts}
\usepackage{amsmath}

\title{\textbf{Analysis on the Real Line HW 2}}
\author{Robert O'Connell}
\date{\today}
\begin{document}
\maketitle
\vspace{5mm}

\subsection*{BS - Page 30 Exercise 18}
Let \(a,b \in \mathbb{R}\), and suppose that for every \(\epsilon > 0\) we have
\(a \leq b + \epsilon\). We want to prove that \(a \leq b\).
\\

To do this we will use a contrapostive proof.

i.e. prove there exists an \(\epsilon\) such that \(a>b \Rightarrow a > b + \epsilon\)
\\

Since \(a>b \Rightarrow a-b > 0\).

Picking epsilon to be \(\frac{a-b}{2}\), which is positive, we see that \(a = b + 2\epsilon\)

Thus \(a > b + \epsilon\)

Which proves that \(a \leq b + \epsilon \Rightarrow a \leq b\)

\subsection*{BS - Page 34 Exercise 5}
We know \(a <x <b\) and \(a<y<b \Rightarrow -b < -y < -a\).

Adding these together we get \(a-b < x-y < b-a\)

Since \(-(a-b) = b-a\), we can use one of the properties of the absolute value function to get \(|x-y| < b-a\).

Interpreting this geometrically, we get that the distance between \(a\) and \(b\) is always greater than the distance between 
\(x\) and \(y\)




\subsection*{BS - Page 38 Exercise 5}
Define \(-S = \{-s | s \in S\}\)

Let \(s_0\) be  \(\inf(S)\)
\begin{itemize}
    \item \(s_0 \geq s \quad \forall \ s \in S\)
    \item If \(z\) is another lower bound for \(S\), i.e. \( z \geq s \quad \forall \ s \in S \), then \(s_0 \geq z \quad \forall \ z\)
\end{itemize}

Using these properties we see 
\begin{itemize}
    \item \(-s_0 \leq -s \quad \forall \ s \in S\)
    \item \( -s_0 \leq z \quad \forall \ z\) such that \( -z \leq -s \quad \forall \ s \in S\)
\end{itemize}

Then
\begin{itemize}
    \item \(-s_0 \leq k \quad \forall \  k \in -S\)
    \item \( -s_0 \leq z \quad \forall \ z\) such that \( -z \leq k \quad \forall \ k \in -S\)
\end{itemize}

Thus we see \(-s_0\) satisfies the properties of being a supremum on the set \(-S\)
\[
-s_0 = \sup (-S) \Rightarrow \inf(S)  = -\sup (-S)
\]



\subsection*{BS - Page 43 Exercise 1}
From class we have that \(\sup(x+A) = x + \sup(A)\), \(A \subset \mathbb{R}, x \in \mathbb{R}\)

i.e. \(\sup(1- \frac{1}{n}| n \in \mathbb{N}) = 1 + \sup(-\frac{1}{n}| n \in \mathbb{N})\)
\\

Now we need to show, \(\sup(-\frac{1}{n}| n \in \mathbb{N}) = 0\)

\begin{itemize}
    \item \(n > 0 \Rightarrow -\frac{1}{n} < 0\) \(\quad \forall \ n \in \mathbb{N}\)
    \\ 
    0 is an upper bound for the set

    \item Let \(s = \sup(-\frac{1}{n}| n \in \mathbb{N})\)
    
    Hence \(s \leq 0\)

    If \(s < 0 \Rightarrow -s > 0\) then by the Archimedean Property 
    \begin{itemize}
        \item \(-sn > 1\) for some \(n \in \mathbb{N}\)
        \item \(-s > \frac{1}{n}\)
        \item \(s < - \frac{1}{n}\)
        \item \(s\) is not an upper bound, a contradiction
    \end{itemize}

    Hence \(s = 0\)
\end{itemize}

\(\sup(1- \frac{1}{n}| n \in \mathbb{N}) = 1 + 0 = 1\)




\end{document}